% Options for packages loaded elsewhere
\PassOptionsToPackage{unicode}{hyperref}
\PassOptionsToPackage{hyphens}{url}
%
\documentclass[
]{article}
\usepackage{amsmath,amssymb}
\usepackage{iftex}
\ifPDFTeX
  \usepackage[T1]{fontenc}
  \usepackage[utf8]{inputenc}
  \usepackage{textcomp} % provide euro and other symbols
\else % if luatex or xetex
  \usepackage{unicode-math} % this also loads fontspec
  \defaultfontfeatures{Scale=MatchLowercase}
  \defaultfontfeatures[\rmfamily]{Ligatures=TeX,Scale=1}
\fi
\usepackage{lmodern}
\ifPDFTeX\else
  % xetex/luatex font selection
\fi
% Use upquote if available, for straight quotes in verbatim environments
\IfFileExists{upquote.sty}{\usepackage{upquote}}{}
\IfFileExists{microtype.sty}{% use microtype if available
  \usepackage[]{microtype}
  \UseMicrotypeSet[protrusion]{basicmath} % disable protrusion for tt fonts
}{}
\makeatletter
\@ifundefined{KOMAClassName}{% if non-KOMA class
  \IfFileExists{parskip.sty}{%
    \usepackage{parskip}
  }{% else
    \setlength{\parindent}{0pt}
    \setlength{\parskip}{6pt plus 2pt minus 1pt}}
}{% if KOMA class
  \KOMAoptions{parskip=half}}
\makeatother
\usepackage{xcolor}
\setlength{\emergencystretch}{3em} % prevent overfull lines
\providecommand{\tightlist}{%
  \setlength{\itemsep}{0pt}\setlength{\parskip}{0pt}}
\setcounter{secnumdepth}{-\maxdimen} % remove section numbering
\ifLuaTeX
  \usepackage{selnolig}  % disable illegal ligatures
\fi
\usepackage{bookmark}
\IfFileExists{xurl.sty}{\usepackage{xurl}}{} % add URL line breaks if available
\urlstyle{same}
\hypersetup{
  hidelinks,
  pdfcreator={LaTeX via pandoc}}

\author{}
\date{}

\begin{document}

\section{The Program Arrives Late}\label{the-program-arrives-late}

\subsection{Theory Lag and Programming Before
Understanding}\label{theory-lag-and-programming-before-understanding}

Watson Hartsoe

\textbf{Final Working Paper - August 2026}

\subsubsection{Abstract}\label{abstract}

Natural-language programming is usually described as a change in
notation: instead of writing formal code, a programmer states an
intention in ordinary language and a generative model translates that
intention into executable form. This account places too much explanatory
weight on the prompt and assumes too much coherence in the intention
that precedes it. Programming from incomplete specifications long
predates large language models, and designers routinely act before
problems are fully determined. The more consequential change introduced
by generative systems is that increasingly heterogeneous forms of
unresolvedness can remain computationally productive. Procedures,
categories, examples, aesthetic criteria, contextual references, and
even the meaning of the desired transformation can remain unsettled
while substantial artifacts are already being generated, inspected,
rejected, and revised.

I call the resulting condition \textbf{theory lag}: the interval in
which a human-computational system can already produce and discriminate
among operative realizations while the human's explicit theory remains
insufficient to explain those discriminations or guide unforeseen
modification. Theory lag reveals three mechanisms that often appear
identical at the interface: \textbf{task location}, in which an
expression elicits behavior already latent in a model;
\textbf{specification completion}, in which provisional realizations
expose or supply commitments missing from the initial description; and
\textbf{tacit steering}, in which human judgment guides generation
without becoming fully articulable as specification. These mechanisms
shift the unit of programming away from the isolated prompt toward a
recursive relation among description, situated interpretation,
realization, judgment, and revision.

Following Peter Naur's account of programming as theory building, I
argue that successful generation is an inadequate test of whether a
program's theory has been preserved. The stronger test is
\textbf{modification competence}: whether a representation can support
principled change when circumstances depart from those under which it
was produced. Natural-language programming therefore raises a question
more difficult than whether prompts count as code: \textbf{what must be
preserved when execution can begin before understanding is complete?}

\textbf{Keywords:} natural-language programming; generative AI;
specification; program synthesis; prompting; situated action; tacit
knowledge; theory building

\subsection{1. Before the Requirement
Exists}\label{before-the-requirement-exists}

Imagine telling a generative system:

\begin{quote}
Make the house remember everyone who has lived here.
\end{quote}

A house appears whose walls display photographs of former residents.

No.

\begin{quote}
The house should remember them. It should not display records about
them.
\end{quote}

The next version alters its lighting according to patterns associated
with former occupants.

Closer.

\begin{quote}
Memory should change what the house does, but nobody living there should
be able to inspect a complete archive.
\end{quote}

A third house appears.

The common description of this exchange is straightforward. The first
prompt was underspecified. The model produced an incorrect
implementation. Subsequent prompts added missing requirements. Through
iteration, the user approached a sufficiently complete specification.

But this description assumes that the missing requirement already
existed.

Perhaps before seeing the photograph-covered walls, the user had never
distinguished \textbf{memory as archive} from \textbf{memory as altered
behavior}. The first realization does not merely violate an existing
specification. It makes a distinction available that the original
description did not contain.

The distinction can then become part of the next instruction.

Something has therefore happened between intention and specification.

The artifact has participated in the production of the requirement by
which the artifact is judged.

This is where natural-language programming becomes difficult to describe
using the conventional sequence:

\texttt{requirement\ -\textgreater{}\ program\ -\textgreater{}\ execution}

In generative interaction, the relation can become recursive:

\texttt{provisional\ description\ -\textgreater{}\ realization\ -\textgreater{}\ judgment\ -\textgreater{}\ revised\ description}

and, more importantly,

\texttt{realization\ -\textgreater{}\ newly\ available\ distinction\ -\textgreater{}\ reconstituted\ requirement}.

The result is not simply a faster route from intention to
implementation. The route can alter the intention itself.

\subsection{2. Incompleteness Is Not
New}\label{incompleteness-is-not-new}

It would be easy to overstate what has changed.

Programming did not previously require every detail of a program to be
settled before computation began. Program synthesis has long accepted
incomplete specifications. Programming-by-example can infer procedures
from demonstrations rather than explicit algorithms. Interactive program
synthesis makes incompleteness central to the programming process, using
incremental computation and feedback to refine what users want (Le et
al.~2017). Designers likewise do not wait until a problem is fully
articulated before acting. Sketches, prototypes, mock-ups, and
provisional arrangements regularly disclose properties of a problem that
could not be derived from reflection alone. Recent empirical work in
visualization design likewise describes expert practitioners making
unresolved situations actionable through provisional local moves
(Parsons et al.~2026).

The historical claim, then, cannot be:

\begin{quote}
Generative AI makes it possible to begin before the specification is
complete.
\end{quote}

That happened already.

A better question is:

\begin{quote}
\textbf{What can remain unresolved while useful computation continues?}
\end{quote}

Earlier systems commonly restrict unresolvedness through an established
representational regime. A synthesis system may leave the procedure
unknown while fixing the domain-specific language in which the solution
must be expressed. A compiler may postpone values until runtime while
preserving a highly explicit operational semantics. A
programming-by-example system may infer transformations while
constraining the kinds of examples it can interpret.

Generative models tolerate a more heterogeneous field of missingness.

The unresolved part can be procedural:

\begin{quote}
Find a better way to organize these records.
\end{quote}

It can be aesthetic:

\begin{quote}
This feels too corporate.
\end{quote}

Deictic:

\begin{quote}
Keep this part.
\end{quote}

Analogical:

\begin{quote}
Make it behave more like a garden than a filing cabinet.
\end{quote}

Narrative:

\begin{quote}
The house should remember, but it should never confess.
\end{quote}

Evaluative:

\begin{quote}
Again. The last one understood the concept too literally.
\end{quote}

These are not merely unspecified parameter values. They are descriptions
whose operative significance depends on contextual interpretation,
learned regularities, examples, prior artifacts, and subsequent
judgment.

The transformation is therefore better described as an expansion in the
\textbf{kinds of unresolvedness that can remain computationally viable}.

The user is permitted to begin before many distinctions have been
translated into the vocabulary from which conventional specifications
are made.

\subsection{3. Theory Lag}\label{theory-lag}

Peter Naur's ``Programming as Theory Building'' provides a useful way to
understand why this matters.

For Naur, the central achievement of programming is not the program text
itself. The programmer develops a theory connecting the program to the
portion of the world it is intended to handle. Possessing that theory
means more than knowing what the existing program does. It includes
knowing why its structures are present, how they correspond to relevant
affairs, and how they ought to change when new demands arise (Naur
1985).

This is why maintenance is so important to Naur's account. Program text
and documentation can survive while the understanding needed to modify
the program intelligently disappears.

Generative programming introduces a peculiar temporal possibility into
this picture.

A functioning realization may appear \textbf{before} the human possesses
anything like the theory Naur describes.

The human can run it.

Reject it.

Notice a difference.

Retain one behavior.

Discard another.

Change the request because the generated artifact reveals that the
original request was badly framed.

Substantial computation can proceed while explicit understanding remains
incomplete.

I call this \textbf{theory lag}.

\begin{quote}
\textbf{Theory lag is the interval in which a human-computational system
can already produce, evaluate, and revise operative realizations while
the human's explicit theory remains insufficient to account for the
distinctions guiding those revisions or to determine how the system
should respond to unforeseen change.}
\end{quote}

Theory lag should not be confused with ignorance. A practitioner inside
the lag may possess enormous expertise. They may know immediately that
something is wrong. They may know which candidate is stronger. They may
recognize a violation before they can state the violated rule.

The important asymmetry is simply this:

\texttt{capacity\ to\ discriminate\ \textgreater{}\ capacity\ to\ explicate}.

Generative systems make this asymmetry computationally productive.

The next substantial realization does not have to wait for the gap to
close.

\subsection{4. When the Artifact Produces the
Requirement}\label{when-the-artifact-produces-the-requirement}

Theory lag matters most where realization does not simply test a
requirement but helps make that requirement interpretable.

Suppose the user asks:

\begin{quote}
Make the interface calm.
\end{quote}

The system produces a pale interface with generous white space.

The user rejects it:

\begin{quote}
Not spa calm. Calm because nothing is competing for attention.
\end{quote}

That correction contains a distinction.

But where was the distinction before the failed interface?

One possibility is that the user already possessed it and merely failed
to articulate it. The generation functions as requirements elicitation.

Another possibility is more interesting. The contrast between \emph{spa
calm} and \emph{noncompetitive calm} becomes available only because a
concrete realization has separated them.

The failed artifact has performed an epistemic operation.

It has given the practitioner something sufficiently definite to negate.

This yields a different account of iteration:

\texttt{D\_t\ -\textgreater{}\ Y\_t\ -\textgreater{}\ J\_t\ -\textgreater{}\ D\_(t+1)}

where \texttt{D\_t} is not a fixed requirement receiving additional
detail. It is a description whose meaning may itself be transformed by
the judgment applied to \texttt{Y\_t}.

The next description can therefore contain distinctions that did not
exist propositionally in the previous one.

That is why \emph{feedback} is too weak a word.

Feedback suggests that a known target exists and that error tells us how
far the current realization lies from it.

Generative iteration can instead change the target space.

The user is not merely moving toward an object.

The succession of objects can change what \emph{toward} means.

\subsection{5. Three Different Things Called
Prompting}\label{three-different-things-called-prompting}

The interface makes this process look simpler than it is.

A person types a small amount of language. A system produces a much
larger structure. It is tempting to explain every such event by saying
that the model ``filled in the gaps.''

At least three different mechanisms can produce this appearance.

\subsubsection{5.1 Task location}\label{task-location}

Reynolds and McDonell propose that some prompting behavior is better
understood as \textbf{task location}. Rather than teaching the model a
new task at inference time, an expression can help locate behavior
already learned during training (Reynolds and McDonell 2021).

Under this account:

\texttt{expression\ -\textgreater{}\ latent\ behavioral\ region\ -\textgreater{}\ realization}.

A sparse prompt can produce elaborate behavior because the prompt is
sparse while the larger operative system is not.

``Translate this as a meticulous philologist'' may work not because the
sentence specifies philological translation but because it selects among
capabilities already available within the model.

\subsubsection{5.2 Specification
completion}\label{specification-completion}

Other cases require stronger interpretive commitment.

``Make the house remember'' does not merely select one obvious task. The
realization must settle dimensions the user has not determined. Memory
becomes behavioral rather than archival; spatial rather than
propositional; persistent rather than episodic.

A generation therefore externalizes one possible completion of an
incomplete description.

That completion can then expose a missing criterion. Empirical research
on prompt programming likewise reports that developers sometimes
discover requirements through interacting with model behavior and
struggle to anticipate which assumptions must be made explicit (Liang et
al.~2024).

\texttt{partial\ description\ -\textgreater{}\ interpretive\ commitment\ -\textgreater{}\ realization\ -\textgreater{}\ criterion\ discovery}.

\subsubsection{5.3 Tacit steering}\label{tacit-steering}

But even specification completion assumes that the process tends toward
an explicit specification.

It may not.

Michael Polanyi's account of tacit knowing begins from the fact that
human competence routinely exceeds what can be stated propositionally.
One can identify a familiar face without possessing an explicit
procedure for recognizing it (Polanyi 2009 {[}1966{]}).

Generative interaction creates an analogous possibility.

A writer can repeatedly reject sentences.

A designer can select one image from eight.

A programmer can say:

\begin{quote}
This architecture is closer, but the state model is still wrong.
\end{quote}

The judgment can be stable and consequential while its underlying
criterion remains only partially articulable.

The process may therefore continue as:

\texttt{candidate\ -\textgreater{}\ discrimination\ -\textgreater{}\ adapted\ candidate\ set\ -\textgreater{}\ discrimination\ -\textgreater{}\ ...}

without terminating in a complete symbolic representation of the
discriminating rule.

Call this \textbf{tacit steering}.

These mechanisms can coexist. But they locate the operative structure
differently.

Task location places much of it in the learned model.

Specification completion distributes it across description and
realization.

Tacit steering leaves an irreducible portion in ongoing human judgment.

A theory of prompting that calls all three ``natural-language
specification'' loses precisely what needs explanation.

\subsection{6. The Prompt Is Not the
Unit}\label{the-prompt-is-not-the-unit}

Consider the prompt:

\begin{quote}
Make this darker.
\end{quote}

As autonomous text, it specifies almost nothing.

\emph{This} may refer to an image, a paragraph, a selected region, a
previously generated object, or an element inside an interactive
environment. \emph{Darker} may concern luminance, color, atmosphere,
rhetoric, or narrative consequence.

Yet within the right situation, the expression may be entirely adequate.

Lucy Suchman's critique of plans helps explain why. A plan can function
as a resource for situated action without determining the detailed
course of that action. The significance of a representation depends
partly upon the circumstances in which it is taken up (Suchman 1987).

This does not mean a language model is situated in the same sense as a
human actor. It means that \textbf{precedence should not be confused
with determination}. An expression occurring before an operation need
not contain the operation in compressed form.

For generative systems, an operative event depends on something closer
to:

\texttt{E\_t\ =\ (p\_t,\ c\_t,\ m,\ a\_t,\ y\_(t-1),\ j\_(t-1))}

where:

\begin{itemize}
\tightlist
\item
  \texttt{p\_t} is the visible expression;
\item
  \texttt{c\_t} is the addressable context;
\item
  \texttt{m} is the interpreting model;
\item
  \texttt{a\_t} is the available repertoire of operations;
\item
  \texttt{y\_(t-1)} is the previous realization;
\item
  \texttt{j\_(t-1)} is prior human judgment.
\end{itemize}

Change one of these while leaving the words untouched and the operative
meaning may change.

The prompt is therefore not a self-sufficient program.

It is a \textbf{move within a changing coupled state}.

This explains the curious power of expressions such as:

\begin{quote}
This.
\end{quote}

\begin{quote}
Again.
\end{quote}

\begin{quote}
Like the last one, except\ldots{}
\end{quote}

\begin{quote}
Keep the left half.
\end{quote}

Their effectiveness comes precisely from information that is not inside
them.

\subsection{7. Description Does Not Become
Operation}\label{description-does-not-become-operation}

This also weakens a seductive description of generative systems:

\begin{quote}
Description becomes operation.
\end{quote}

A description does not acquire an intrinsic operational essence.

Consider:

\begin{quote}
Delete this.
\end{quote}

In a novel, nothing happens.

In an ordinary chat, the recipient may ask what \emph{this} refers to.

In an image editor, a selected object may disappear.

In an agentic environment, a tool invocation may be produced.

In another environment, the operation may be blocked because the user
lacks permission.

The words are constant.

The operative arrangement changes.

J. L. Austin's account of performative utterances is useful here because
the familiar slogan ``words do things'' is only half of his argument.
Performative force depends upon circumstances, procedures, roles, and
conditions under which an act succeeds or misfires (Austin 1975).

Computational systems have analogous conditions.

An expression becomes consequential only when relevant conditions hold:

\texttt{F\ =\ \{resolvable\ referent,\ recognized\ operation,\ authorized\ actor,\ available\ capability,\ admissible\ state,\ successful\ transition\}}.

The important phenomenon is therefore not executable description in
isolation.

It is \textbf{interpretive coupling}: arrangements in which an
interpreted representation can modify what states become reachable next.

This formulation returns the machinery that ``description becomes
operation'' makes disappear.

It also exposes power.

Who defines what can be referenced?

Who chooses the interpreter?

Who grants tool access?

Which state transitions are available?

Which utterances are treated as instructions?

The apparent power of language is partly the power of the infrastructure
to which language has been connected.

\subsection{8. The Program Arrives
Late}\label{the-program-arrives-late-1}

Return to the remembering house.

After enough iterations, the user may eventually be able to write
something much more explicit:

\begin{itemize}
\tightlist
\item
  memory changes behavior rather than displaying records;
\item
  traces of past occupants alter present affordances;
\item
  memory is partial rather than total;
\item
  occupants cannot inspect it directly;
\item
  forgetting remains possible;
\item
  remembrance should be sensed indirectly.
\end{itemize}

This object now looks recognizably like a specification.

But it did not exist at the beginning.

It is partly a retrospective stabilization of distinctions produced
during interaction.

The theory arrived after versions of the system had already been
generated.

Hence the central temporal reversal:

\begin{quote}
\textbf{The program, understood as an explicit account of what must
remain true, can arrive after the system has begun to run.}
\end{quote}

This does not mean there was no prior intention.

Nor does it mean the model authored the theory.

It means that the relation between theory and realization can become
reciprocal.

At iteration \texttt{t}:

\texttt{T\_t\ -\textgreater{}\ D\_t\ -\textgreater{}\ Y\_t\ -\textgreater{}\ J\_t\ -\textgreater{}\ T\_(t+1)}.

\texttt{T\_(t+1)} need not simply contain more details than
\texttt{T\_t}.

A generated realization may reveal that the relevant distinction was
wrong, that two requirements were incompatible, or that the object under
construction has become something different from the one initially
imagined.

The theory can change kind.

\subsection{9. Why Regeneration Proves Too
Little}\label{why-regeneration-proves-too-little}

Once a sufficiently explicit theory has emerged, another temptation
appears.

Perhaps it can be compressed.

A long program, essay, world, or dissertation could be represented by a
smaller generative object capable of reproducing recognizable
realizations.

If an 80,000-word argument can be regenerated from a compact prompt
system, perhaps the prompt system is the deeper intellectual artifact.

Naur gives us a reason to resist this conclusion.

A representation can reproduce an artifact without preserving the theory
required to alter it intelligently.

Let \texttt{R(P)} denote the fidelity with which representation
\texttt{P} regenerates previously observed realizations.

Now define \texttt{M(P,\ ΔW)} as the competence with which \texttt{P},
together with its interpreter, supports modification under an unforeseen
change in relevant world conditions \texttt{ΔW}.

Nothing guarantees:

\texttt{R(P)\ increases\ -\textgreater{}\ M(P,\ ΔW)\ increases}.

A prompt can reproduce an essay because a model recognizes the genre and
argument pattern.

It can reproduce software because common architecture is supplied by
learned priors.

It can reconstruct an interface because examples constrain imitation.

None of these establishes that the representation preserves why the
existing structures matter.

The harder test is \textbf{modification competence}.

Change the world.

Remove a premise.

Introduce evidence the argument cannot absorb.

Require the remembering house to forget.

Change the jurisdiction governing the software.

Move the interface from a single user to collaborative ownership.

Now ask:

Which structures should change?

Which should remain?

Why?

A representation preserving the relevant program theory should support
coherent modification under such counterfactual pressure.

Regeneration asks whether the old world can be rebuilt.

Modification asks whether the representation knows how to survive a new
one.

\subsection{10. Making Theory Lag
Falsifiable}\label{making-theory-lag-falsifiable}

Theory lag should not survive merely because it is an attractive
description.

It suggests at least three empirical tests.

\subsubsection{Criterion emergence}\label{criterion-emergence}

Before generation, elicit every criterion a participant can state.

After each realization, record new criteria.

Distinguish criteria that were:

\begin{itemize}
\tightlist
\item
  consciously known but omitted;
\item
  recognized as relevant only after a violation;
\item
  created by an unexpected possibility;
\item
  merely post-hoc explanations.
\end{itemize}

If generations regularly precede stable articulation of consequential
criteria, theory lag gains empirical support.

\subsubsection{Judgment before
explanation}\label{judgment-before-explanation}

Present alternatives and require immediate ranking or rejection before
asking participants to explain their judgment.

Then compare consistency of discrimination with consistency and
predictive power of explanation.

If people reliably discriminate among candidates while their explicit
rules remain unstable or incomplete, operative judgment is outrunning
explicit theory.

\subsubsection{Regeneration versus
modification}\label{regeneration-versus-modification}

Construct a compact generative representation that reliably reproduces
an existing artifact.

Then introduce unforeseen requirements.

Measure separately:

\begin{itemize}
\tightlist
\item
  regeneration fidelity;
\item
  explanation of existing structures;
\item
  identification of relevant invariants;
\item
  quality of modification;
\item
  preservation of original rationale.
\end{itemize}

If regeneration succeeds while principled modification collapses,
realization has survived without sufficient theory.

These tests could also kill the concept.

If criteria are generally explicit before generation, theory lag adds
little beyond ordinary requirements elicitation.

If post-generation explanations are mostly unstable rationalizations,
realization may be producing preference drift rather than theory.

If modification competence closely tracks regeneration fidelity, the
distinction between reproducing a system and preserving its theory may
be weaker than proposed.

And if theory lag proves no greater in generative practice than in
ordinary skilled design, its significance will need to be narrowed
further.

\subsection{11. The Economy of Provisional
Worlds}\label{the-economy-of-provisional-worlds}

If incomplete specification, situated action, iteration, and tacit
judgment all predate generative AI, what remains distinctive?

Perhaps not a new logical structure.

Perhaps a new \textbf{economy of provisional realization}.

A sketch takes time.

A prototype takes more.

A functioning interface, substantial essay, architectural visualization,
or alternative codebase traditionally requires increasing quantities of
labor before it can be judged.

Generative systems compress that cost.

An unresolved description can produce a substantial candidate quickly
enough that candidate-making itself becomes a routine mode of thought.

The number and heterogeneity of realizations available before a
practitioner must commit can rise dramatically.

Call this \textbf{iteration bandwidth}.

High iteration bandwidth can improve inquiry.

More realizations expose more differences.

Previously invisible criteria become visible.

Unexpected possibilities alter the problem.

But the same abundance produces the opposite danger.

Every difficult passage now has another version.

Every unresolved interface has another arrangement.

Every architectural contradiction has another plausible image.

The machine may allow the practitioner to route around a difficulty
before the difficulty becomes intellectually expensive enough to
understand.

The danger is therefore not simply automation.

It is \textbf{evasion through abundance}.

Slow work sometimes forces theory because the cost of continuing without
understanding becomes intolerable.

A system capable of generating indefinitely plausible next moves can
keep theory lag open.

\subsection{12. Designing for the Gap}\label{designing-for-the-gap}

This suggests a different agenda for natural-language programming
environments.

Most systems are optimized to reduce friction between intention and
result.

Better inference.

Shorter prompts.

More automatic context.

Fewer clarifying questions.

Greater ability to infer what the user ``really meant.''

But if theory lag is real, eliminating every friction can conceal the
very places where theory remains inadequate.

A programming environment built around theory lag might therefore do
something unusual.

It might preserve repeated corrections and ask whether they imply a new
invariant.

It might distinguish a local patch from a change in the governing
problem.

It might retain rejected realizations because they reveal boundaries of
the emerging theory.

It might show which behaviors come primarily from the prompt and which
appear to be supplied by model priors.

It might periodically introduce controlled perturbations:

What if this assumption disappears?

What if the data model changes?

What if there are three owners rather than one?

What if memory must decay?

The point would not be to force every intuition into a formal
specification.

Some judgment may remain tacit.

The point would be to expose the difference between:

\textbf{a system that continues to produce acceptable outputs}

and

\textbf{a practitioner who increasingly understands what must remain
true.}

These are not the same accomplishment.

\subsection{13. Programming as Theory Building After
Execution}\label{programming-as-theory-building-after-execution}

Naur placed theory building at the center of programming because program
text alone could not preserve the understanding required for intelligent
change.

Generative systems do not make that theory unnecessary.

They alter when it may be acquired.

The practitioner can now encounter substantial realizations while
important parts of the theory remain unsettled.

They can discover criteria through failure.

Locate behaviors through sparse expressions.

Steer by judgments they cannot fully articulate.

Allow previously informal descriptions to participate in consequential
state changes.

Programming therefore becomes neither pure specification nor pure
selection.

It becomes a process in which \textbf{theory and realization can develop
together}.

That is the stronger interpretation of the familiar loop:

\texttt{describe\ -\textgreater{}\ generate\ -\textgreater{}\ inspect\ -\textgreater{}\ correct\ -\textgreater{}\ regenerate}.

The loop is not important because iteration improves prompts.

It is important because each traversal can change the theory under which
the artifact is judged.

The real output of the loop is therefore not only the latest artifact.

It is also the evolving account of what the artifact is supposed to be.

\subsection{Conclusion}\label{conclusion}

Natural-language programming should not be understood primarily as
programming in English.

Nor does its significance lie in proving that prompts are a new species
of code.

Programming from incomplete specifications already existed. Skilled
action already exceeded explicit rules. Designers already learned from
prototypes. Plans already depended on situations. Language already
acquired consequences inside institutional arrangements.

Generative systems reorganize these older relations by making
\textbf{heterogeneous unresolvedness cheaply operable}.

A programmer can now submit an intention before its consequential
distinctions have been completely articulated, encounter substantial
realizations, recognize differences that were previously unavailable,
and feed those judgments back into the process without first translating
them into a formal language.

This produces theory lag: execution and evaluation can proceed ahead of
explicit understanding.

Inside that lag, the apparent simplicity of prompting conceals different
mechanisms. Sometimes a small expression locates a capability already
present in the model. Sometimes a realization supplies commitments that
help constitute a specification. Sometimes the human steers through
judgments that never become completely explicit.

None of these makes theory obsolete.

They make its delayed arrival possible.

And delayed theory creates a problem that successful generation cannot
solve.

A representation may regenerate an artifact without preserving the
understanding required to modify it coherently when circumstances
change. The stronger test is therefore not reproduction but modification
competence.

The old programming question was:

\begin{quote}
\textbf{How precisely must the program be specified before it can run?}
\end{quote}

Generative systems make another question unavoidable:

\begin{quote}
\textbf{How far can a program run before its programmer understands what
must remain true?}
\end{quote}

That is not merely a question about prompts.

It is a question about the temporal order of programming itself.

When realization can precede understanding, the central task becomes
preserving enough friction, history, judgment, and theory that
understanding can still catch up.

The program can arrive late.

It cannot be permitted never to arrive.

\subsection{References}\label{references}

Austin, J. L. 1975. \emph{How to Do Things with Words}. 2nd ed.~Edited
by J. O. Urmson and Marina Sbisà. Cambridge, MA: Harvard University
Press. First published 1962.

Le, Vu, Daniel Perelman, Oleksandr Polozov, Mohammad Raza, Abhishek
Udupa, and Sumit Gulwani. 2017. ``Interactive Program Synthesis.''
arXiv:1703.03539.

Liang, Jenny T., Melissa Lin, Nikitha Rao, and Brad A. Myers. 2024.
``Prompts Are Programs Too! Understanding How Developers Build Software
Containing Prompts.'' arXiv:2409.12447.

Naur, Peter. 1985. ``Programming as Theory Building.''
\emph{Microprocessing and Microprogramming} 15 (5): 253-261.
https://doi.org/10.1016/0165-6074(85)90032-8.

Parsons, Paul C., Prakash Shukla, Phuong Bui, Srishti Agrawal, and Ali
Baigelenov. 2026. ``Situatedness in Visualization Design: Making
Unresolved Work Actionable.'' arXiv:2608.08274.

Polanyi, Michael. 2009 {[}1966{]}. \emph{The Tacit Dimension}. Chicago:
University of Chicago Press.

Reynolds, Laria, and Kyle McDonell. 2021. ``Prompt Programming for Large
Language Models: Beyond the Few-Shot Paradigm.'' arXiv:2102.07350.

Suchman, Lucy A. 1987. \emph{Plans and Situated Actions: The Problem of
Human-Machine Communication}. Cambridge: Cambridge University Press.

\end{document}
